% Part: sets-functions-relations
% Chapter: sets
% Section: important-sets

\documentclass[../../../include/open-logic-section]{subfiles}

\begin{document}

\olfileid[id]{sfr}{set}{imp}
\olsection{Beberapa Himpunan Penting}

\begin{ex}
Kita terutama akan membahas himpunan dengan objek matematika sebagai
!!{element}s. Empat himpunan semacam itu cukup penting sehingga
mempunyai nama khusus:
\begin{multline*}
    \Nat = \{0, 1, 2, 3, \ldots\} \\
    \shoveright{\text{himpunan bilangan asli (dengan 0)}}\\
    \shoveleft{\Int = \{\ldots, -2, -1, 0, 1, 2, \ldots\}} \\
    \shoveright{\text{himpunan bilangan bulat}}\\
    \shoveleft{\Rat = \Setabs{\nicefrac{m}{n}}{m, n \in \Int\text{ dan }n \neq 0}}\\
    \shoveright{\text{himpunan bilangan rasional}}\\
    \shoveleft{\Real = (-\infty, \infty)}\\
    \text{himpunan bilangan real (kontinuum)}
\end{multline*}
Dalam teks ini, konvensi bilangan asli mencakup nol, sebagaimana
dinyatakan oleh definisi di atas.
Semua ini merupakan himpunan \emph{tak hingga}, yakni masing-masing
mempunyai !!{element}s yang banyaknya tak hingga.

Ketika bergerak melalui himpunan-himpunan ini, kita menambahkan
\emph{lebih banyak} bilangan ke dalam persediaan kita. Jelas bahwa
$\Nat \subseteq \Int \subseteq \Rat \subseteq \Real$: setiap bilangan
asli merupakan bilangan bulat; setiap bilangan bulat merupakan bilangan
rasional; dan setiap bilangan rasional merupakan bilangan real. Jelas
pula bahwa $\Nat \subsetneq \Int \subsetneq \Rat$, sebab $-1$ adalah
bilangan bulat tetapi bukan bilangan asli, sedangkan $\nicefrac{1}{2}$
adalah bilangan rasional tetapi bukan bilangan bulat. Tidak begitu
jelas bahwa $\Rat \subsetneq \Real$, yaitu bahwa ada bilangan real yang
bukan bilangan rasional\oliflabeldef{sfr:arith:real:realline}{, tetapi
kita akan kembali membahasnya dalam
\olref[arith][real]{realline}}{}.

Kadang-kadang kita juga akan menggunakan himpunan bilangan bulat
positif $\PosInt = \{1, 2, 3, \dots\}$ dan himpunan yang hanya memuat
dua anggota pertama himpunan bilangan asli menurut konvensi di atas,
yaitu $\Bin = \{0, 1\}$.
\end{ex}

\begin{tagblock}{compsci}
\begin{ex}[Untai]
Contoh menarik lainnya adalah himpunan $A^{*}$ dari \emph{untai
berhingga} atas alfabet $A$: setiap barisan berhingga anggota dari
$A$ merupakan untai atas $A$. Kita menyertakan \emph{untai kosong
$\Lambda$} di antara untai-untai atas~$A$, untuk setiap alfabet~$A$.
Sebagai contoh,
\begin{multline*}
\Bin^*
=\{\Lambda,0,1,00,01,10,11,\\
000,001,010,011,100,101,110,111,0000,\ldots\}.
\end{multline*}
Jika $x=x_{1}\ldots x_{n}\in A^{*}$ adalah untai yang terdiri atas $n$
``huruf'' dari $A$, kita mengatakan bahwa \emph{panjang} untai itu
adalah~$n$ dan menulis $\len{x}=n$.
\end{ex}
\end{tagblock}

\begin{ex}[Barisan tak hingga]
Untuk setiap himpunan $A$, kita juga dapat meninjau himpunan~$A^\omega$
dari barisan tak hingga !!{element}s dari~$A$. Barisan tak hingga
$a_1a_2a_3a_4\dots$ terdiri atas daftar objek satu arah yang tak
berhingga, dan setiap objek dalam daftar itu merupakan !!a{element}
dari~$A$.
\end{ex}

\end{document}
